\documentclass[10pt,a4paper]{amsart}
\usepackage[margin=19mm]{geometry}
\usepackage{amsmath,amssymb,mathtools}
\usepackage[scr=rsfs,cal=euler]{mathalfa}
\usepackage{xfrac}
\usepackage[unicode,hidelinks,bookmarks=false]{hyperref}
\newcommand{\ie}{i.e.\ }
\newcommand{\cf}{cf.\ }
\newcommand{\resp}{respectively\ }
\newcommand{\Subsection}{\subsection}
\providecommand{\nobreakdash}{\nobreak\hbox{-}}
\setlength{\emergencystretch}{2em}
\multlinegap0pt
\usepackage[normalem]{ulem}
\usepackage{amssymb}
\usepackage{xcolor}
\newtheorem{proposition}{Proposition}
\newcommand {\N} {\mathbb{N}}
\newcommand{\conj}[2]{\left \{ {#1} \, : \, {#2} \right \}}
\newcommand{\norma}[1]{\| #1 \|}
\title{A $p$-summability approach to the weak maximizing property}
\author{Oscar I. Blanco A. \, and Vinícius C. C. Miranda}
\date{Source: arXiv:2609.03988v1 (CC BY 4.0)}
\begin{document}
\maketitle

\noindent\emph{Source and licence.} A $p$-summability approach to the weak maximizing property. Version arXiv:2609.03988v1 (CC BY 4.0), distributed by arXiv under the Creative Commons Attribution 4.0 International licence, http://creativecommons.org/licenses/by/4.0/, as recorded in the arXiv licence metadata for this exact version and shown on the abstract page https://arxiv.org/abs/2609.03988v1. Full licence notice: \texttt{licenses/arxiv-2609.03988-CC-BY-4.0.txt}.\par\smallskip

\noindent\textit{Editorial scope.} Counted unit: the article's proposition primeiro (source0.tex lines 87--93). The article's complete original proof of it is delivered with it (source0.tex lines 95--112, one \texttt{proof} environment of 18 lines). No mathematics is written, repaired or re-derived: the statement and the proof are the article's own lines, in the article's own notation, and no retained line is substituted or re-flowed.  No further source-local context is needed: every cross-reference of every retained line resolves inside the two delivered spans.  Nothing was omitted from any retained span, so the package carries no omission record.  No retained line contains a citation, so the wrapper carries no bibliography and no bibliography entry.  No retained line contains a figure environment, an image-inclusion command or drawing code, so no image is delivered and the package declares no asset dependency.  Editorial formatting only, in addition: an AMS \texttt{amsart} preamble that restates the article's own declarations and macros used by the retained lines, namely the environments \texttt{proposition} on separate counters, together with the standard packages those lines need.  The retained environments print in this document as Proposition 1 in order of appearance, whereas the article prints the same items with the numbering of its own full document.  The complete native source of the article as distributed in the arXiv e-print for this exact version is bundled unchanged as \texttt{sources/209332/source0.tex}.
\begin{proposition}\label{primeiro}
Let $X$ and $Y$ be Banach spaces, let $T\in\mathcal L(X;Y)$, and let
$1\leq p<\infty$. Suppose that $T$ admits a maximizing sequence
$(x_n)_n\subset S_X$ which is not weakly $p$-summable. Then, for every
$1\leq q<\infty$, $T$ admits a maximizing sequence which is
not weakly $q$-summable.
\end{proposition}

\begin{proof}
Since $(x_n)_n\notin \ell_p^w(X)$, there exists $x^*\in X^*$ such that $(x^*(x_n))_n\notin \ell_p.$
In particular, the set
$A:=\conj{n \in \N}{x^*(x_n) \neq 0}$
is infinite. Fix $1\leq q<\infty$. We construct the sequence $(m_n)_n \subset \N$ as follows: for each $n\in A$, choose $m_n\in\N$ such that  $m_n |x^*(x_n)|^q\geq 1$. If $n\notin A$, choose $m_n=1$. Define
$$ (z_k)_k   =  (\underbrace{x_1,\ldots,x_1}_{m_1\text{ times}},   \underbrace{x_2,\ldots,x_2}_{m_2\text{ times}},  \ldots,  \underbrace{x_n,\ldots,x_n}_{m_n\text{ times}},    \ldots) \subset S_X.$$
We claim that $(z_k)_k$ is a maximizing sequence for $T$. Indeed, since
$(x_n)_n$ is a maximizing sequence for $T$, given $\varepsilon > 0$, there exists $n_0 \in \N$ such that $\norma{Tx_n} > \norma{T} - \varepsilon$ for every $n \geq n_0$. Taking $k_0 = m_1 + \cdots + m_{n_0 -1} + 1$, by the definition of $(z_k)_k$, we have that for each $k \geq k_0$, $z_k = x_n$ for some $n \geq n_0$. Hence $\norma{Tz_k} > \norma{T} - \varepsilon$ for every $k \geq k_0$, proving that $\displaystyle \lim_{k \to \infty} \norma{Tz_k} = \norma{T}$.
On the other hand,
$$    \sum_{k=1}^{\infty} |x^*(z_k)|^q
    =
    \sum_{n=1}^{\infty} m_n |x^*(x_n)|^q
    \geq
    \sum_{n\in A} 1
    =
    \infty.$$
Thus $(x^*(z_k))_k\notin \ell_q$, and consequently,  $(z_k)_k\notin \ell_q^w(X).$
\end{proof}

\end{document}
